\documentclass{article}
\usepackage[utf8]{inputenc}
\usepackage{amsmath, amssymb, amsthm}
\usepackage{geometry}
\geometry{a4paper, margin=1in}

\begin{document}

\title{Formulação Matemática do Modelo SDDP}
\author{Extração Automática de Código}
\date{}
\maketitle

\section{Conjuntos}
\begin{itemize}
    \item $\mathcal{T}$: Conjunto de estágios (meses) do horizonte de planejamento.
    % Definido em: meses = todos_meses[1:min(config.num_meses, length(todos_meses))]
    
    \item $\mathcal{S}$: Conjunto de submercados de energia (ex: SE, S, NE, N).
    % Definido em: submercados = unique(data.cenarios.submercado)
    
    \item $\mathcal{U}$: Conjunto de usinas geradoras.
    % Definido em: usinas = unique(data.geracao.usina_cod)
    
    \item $\mathcal{I}_{t,s}$: Conjunto de trades disponíveis no estágio $t \in \mathcal{T}$ e submercado $s \in \mathcal{S}$.
    % Definido em: trades_sub = findall(row -> row.submercado == sub, eachrow(dados.trades))
    
    \item $\Omega_t$: Conjunto de cenários estocásticos amostrados no estágio $t \in \mathcal{T}$.
    % Definido em: cenarios_selecionados = sort(randperm(num_cenarios_total)[1:min(config.num_cenarios, num_cenarios_total)])
\end{itemize}

\section{Parâmetros}
\begin{itemize}
    \item $E_{scale}$: Fator de escala financeira (igual a $10^6$).
    % Definido em: ESCALA = 1e6
    
    \item $L^{cred}$: Limite de crédito ou nível de ruína permitido.
    % Definido em: limite_credito_escala = -100.0  # -100 milhões de reais
    
    \item $H_t$: Número de horas no mês $t$.
    % Definido em: horas_mes(d::Date) = daysinmonth(d) * 24 e dados.horas
    
    \item $V^{leg, B}_{t,s}$: Volume consolidado de contratos de compra existentes no estágio $t$, submercado $s$.
    % Definido em: vol_compra_exist[sub] = nrow(compras) > 0 ? sum(compras.volume_mwm) : 0.0
    
    \item $V^{leg, S}_{t,s}$: Volume consolidado de contratos de venda existentes no estágio $t$, submercado $s$.
    % Definido em: vol_venda_exist[sub] = nrow(vendas) > 0 ? sum(vendas.volume_mwm) : 0.0
    
    \item $P^{leg, B}_{t,s}$: Preço médio dos contratos de compra existentes no estágio $t$, submercado $s$.
    % Definido em: preco_compra_exist[sub] = ... sum(compras.volume_mwm .* compras.preco_r_mwh) / sum(...)
    
    \item $P^{leg, S}_{t,s}$: Preço médio dos contratos de venda existentes no estágio $t$, submercado $s$.
    % Definido em: preco_venda_exist[sub] = ... sum(vendas.volume_mwm .* vendas.preco_r_mwh) / sum(...)
    
    \item $\bar{V}^{B}_{i}$: Limite máximo de volume para o trade de compra $i$.
    % Definido em: dados.trades.limite_compra[i]
    
    \item $\bar{V}^{S}_{i}$: Limite máximo de volume para o trade de venda $i$.
    % Definido em: dados.trades.limite_venda[i]
    
    \item $P^{B}_{i}$: Preço unitário do trade de compra $i$.
    % Definido em: dados.trades.preco_compra[i]
    
    \item $P^{S}_{i}$: Preço unitário do trade de venda $i$.
    % Definido em: dados.trades.preco_venda[i]
    
    \item $G_{t,s}(\omega)$: Geração física agregada no submercado $s$, estágio $t$, cenário $\omega$.
    % Definido em: geracao = Dict(sub => sum(get(dict_geracao, (mes, u, cenario), 0.0) ...
    
    \item $PLD_{t,s}(\omega)$: Preço de Liquidação das Diferenças no submercado $s$, estágio $t$, cenário $\omega$.
    % Definido em: pld = Dict(sub => get(dict_pld, (mes, sub, cenario), 0.0) for sub in submercados)
    
    \item $p(\omega)$: Probabilidade de ocorrência do cenário $\omega$.
    % Definido em: prob = 1.0 / length(cenarios_selecionados)
\end{itemize}

\section{Variáveis de Decisão}
As variáveis de decisão são definidas localmente para o subproblema do estágio $t$:
\begin{itemize}
    \item $x_t \in \mathbb{R}$: Variável de estado representando o saldo de caixa ao final do estágio $t$.
    % Definido em: @variable(sp, caixa, SDDP.State, initial_value=0.0)
    % Nota: Na notação do código, caixa.in = x_{t-1} e caixa.out = x_t
    
    \item $v^{B}_{i,t} \in \mathbb{R}$: Volume negociado no trade de compra $i \in \mathcal{I}_{t,s}$ no estágio $t$. Domínio: $0 \leq v^{B}_{i,t} \leq \bar{V}^{B}_{i}$.
    % Definido em: @variable(sp, 0 <= volume_compra[i=1:num_trades] <= dados.trades.limite_compra[i])
    
    \item $v^{S}_{i,t} \in \mathbb{R}$: Volume negociado no trade de venda $i \in \mathcal{I}_{t,s}$ no estágio $t$. Domínio: $0 \leq v^{S}_{i,t} \leq \bar{V}^{S}_{i}$.
    % Definido em: @variable(sp, 0 <= volume_venda[i=1:num_trades] <= dados.trades.limite_venda[i])
    
    \item $y^{spot}_{s,t} \in \mathbb{R}$: Variável auxiliar (molde estocástico) capturando o resultado financeiro no mercado de curto prazo (spot) no submercado $s$.
    % Definido em: @variable(sp, spot_profit[dados.submercados])
\end{itemize}

\section{Relações entre Variáveis}
\begin{itemize}
    \item $R^{fixo}_t$: Fluxo de caixa determinístico resultante dos contratos legados no estágio $t$.
    \begin{align}
        R^{fixo}_t = \sum_{s \in \mathcal{S}} \frac{\left( P^{leg, S}_{t,s} V^{leg, S}_{t,s} H_t \right) - \left( P^{leg, B}_{t,s} V^{leg, B}_{t,s} H_t \right)}{E_{scale}}
    \end{align}
    % Definido em: lucro_fixo += (receita_venda_exist - custo_compra_exist) / ESCALA
    
    \item $R^{trades}_t$: Fluxo de caixa determinístico total, englobando legado e novos trades do estágio $t$.
    \begin{align}
        R^{trades}_t = R^{fixo}_t + \sum_{s \in \mathcal{S}} \sum_{i \in \mathcal{I}_{t,s}} \frac{\left( v^{S}_{i,t} P^{S}_{i} - v^{B}_{i,t} P^{B}_{i} \right) H_t}{E_{scale}}
    \end{align}
    % Definido em: lucro_trades += (volume_venda[i] * dados.trades.preco_venda[i] - volume_compra[i] * dados.trades.preco_compra[i]) * dados.horas / ESCALA
    
    \item $E_{t,s}(\omega)$: Exposição base ao PLD no cenário $\omega$.
    \begin{align}
        E_{t,s}(\omega) = G_{t,s}(\omega) + V^{leg, B}_{t,s} - V^{leg, S}_{t,s}
    \end{align}
    % Definido em: exposicao_base = ω.geracao[sub] + dados.vol_compra_exist[sub] - dados.vol_venda_exist[sub]
\end{itemize}

\section{Formulação Matemática (Subproblema do Estágio $t$)}

\subsection{Função Objetivo do Estágio}
Maximizar o lucro total do estágio $t$ (recurso imediato) mais a função de custo futuro (Benders cuts implícitos gerados pelo SDDP). O recurso imediato no cenário $\omega$ é modelado como:
\begin{align}
    \max \quad Z_t(\omega) = R^{trades}_t + \sum_{s \in \mathcal{S}} y^{spot}_{s,t}
\end{align}
% Definido em: @stageobjective(sp, lucro_trades + sum(spot_profit[sub] for sub in dados.submercados))

\subsection{Restrições}
\begin{enumerate}
    \item \textbf{Condição Inicial:}
    O estado inicial do fluxo de caixa no momento imediatamente anterior ao primeiro estágio ($t=0$) é nulo:
    \begin{align}
        x_0 = 0
    \end{align}
    % Definido em: @variable(sp, caixa, SDDP.State, initial_value=0.0)
    
    \item \textbf{Transição de Estado (Caixa Acumulado):}
    \begin{align}
        x_t = x_{t-1} + R^{trades}_t + \sum_{s \in \mathcal{S}} y^{spot}_{s,t}
    \end{align}
    % Definido em: @constraint(sp, transicao_caixa, caixa.out == caixa.in + lucro_trades + sum(spot_profit[sub] for sub in dados.submercados))
    
    \item \textbf{Limite de Ruína (Crédito):}
    \begin{align}
        x_t \geq L^{cred}
    \end{align}
    % Definido em: @constraint(sp, limite_ruina, caixa.out >= limite_credito_escala)
    
    \item \textbf{Liquidação no Mercado Spot (Molde Estocástico):}
    A relação entre os volumes negociados, a exposição física e o PLD estocástico é construída através de injeção de coeficientes no JuMP:
    \begin{align}
        y^{spot}_{s,t} - \sum_{i \in \mathcal{I}_{t,s}} v^{B}_{i,t} \left( \frac{PLD_{t,s}(\omega) H_t}{E_{scale}} \right) + \sum_{i \in \mathcal{I}_{t,s}} v^{S}_{i,t} \left( \frac{PLD_{t,s}(\omega) H_t}{E_{scale}} \right) = E_{t,s}(\omega) \left( \frac{PLD_{t,s}(\omega) H_t}{E_{scale}} \right) \quad \forall s \in \mathcal{S}
    \end{align}
    % Definido em: @constraint(sp, spot_profit_eq[sub in dados.submercados], spot_profit[sub] == 0.0)
    % Modificado por:
    % JuMP.set_normalized_rhs(spot_profit_eq[sub], rhs_val)
    % JuMP.set_normalized_coefficient(spot_profit_eq[sub], volume_compra[i], -pld_horas)
    % JuMP.set_normalized_coefficient(spot_profit_eq[sub], volume_venda[i], pld_horas)
\end{enumerate}

\section{Modelagem de: Lucro e Risco}

\subsection{Lucro}
O lucro total de um dado estágio $t$ no cenário $\omega$ é composto pela soma da liquidação das obrigações forward e do balanço financeiro do spot. Ele compõe o objetivo do estágio:
\begin{align}
    Lucro_t(\omega) = R^{trades}_t + \sum_{s \in \mathcal{S}} y^{spot}_{s,t}(\omega)
\end{align}
Onde $y^{spot}_{s,t}(\omega)$ isolando da restrição 3 equivale a:
\begin{align}
    y^{spot}_{s,t}(\omega) = \left[ E_{t,s}(\omega) + \sum_{i \in \mathcal{I}_{t,s}} v^{B}_{i,t} - \sum_{i \in \mathcal{I}_{t,s}} v^{S}_{i,t} \right] \times \frac{PLD_{t,s}(\omega) H_t}{E_{scale}}
\end{align}
% Definido em: @stageobjective(sp, lucro_trades + sum(spot_profit[sub] for sub in dados.submercados))

\subsection{Risco}
A aversão ao risco global do problema ao longo da árvore de cenários é regida por uma combinação convexa entre a Esperança e o Valor em Risco Condicional Médio (AVaR/CVaR), definindo a medida de risco $\rho$ que será otimizada pela política do grafo:
\begin{align}
    \rho[Z] = (1 - \lambda) \cdot \mathbb{E}[Z] + \lambda \cdot \text{AVaR}_{\alpha}[Z]
\end{align}
% Definido em: risk_measure = (1 - config.lambda) * SDDP.Expectation() + config.lambda * SDDP.AVaR(config.alpha)
% E aplicado em: SDDP.train(model, ..., risk_measure=risk_measure)

\section{Regras Obrigatórias e Parâmetros Estruturais}
\begin{itemize}
    \item Nível de confiança do CVaR ($\alpha$): $0.95$.
    % Definido em: 0.95,    # Alpha do CVaR (na função load_sddp_config)
    \item Peso do risco ($\lambda$): $0.01$.
    % Definido em: 0.01,    # Lambda (peso do risco)
    \item Cenários de amostragem por nó (Forward pass): $\min(1000, 2000)$.
    % Definido em: 1000,      # 10 cenários aleatórios (na função load_sddp_config, parâmetro config.num_cenarios)
    \item Simulações Monte Carlo (Avaliação da política): $100$.
    % Definido em: 100      # Simulações (na função load_sddp_config)
    \item Limites de Domínio das Variáveis: Variáveis de estado não possuem limite explícito na declaração (são restritas pelas inequações), e variáveis de trade são sempre estritamente positivas e limitadas superiormente.
    % Definido em: 0 <= volume_compra[i=1:num_trades] <= dados.trades.limite_compra[i]
\end{itemize}

\end{document}